\chapter{El primer teorema, y la mitad que sale gratis}
\label{cap:godel-i}

Con el punto fijo demostrado, el primer teorema de Gödel ocupa seis líneas de Lean y no necesita
ninguna idea nueva. Este capítulo las lee despacio, porque en ellas se ve exactamente
\emph{qué} se usa — y, sobre todo, qué no.

Y luego dice lo que casi ningún libro dice en el mismo sitio: que de las dos mitades del teorema,
en esta formalización sólo una está cerrada.

\section{El argumento}

Supóngase que $G$ \emph{sí} fuera demostrable. Entonces:

\begin{matematica}[Cuatro pasos y una contradicción]
\begin{enumerate}
\item Si $G$ tiene una demostración finitaria, la teoría lo \emph{sabe}: por D1,
  $\text{\ident{axioms}} \derives \Prov(\codigo{G})$. Eso es exactamente lo que el
  capítulo~\ref{cap:d1} demostró.
\item De $\Prf G$ se pasa a $\text{\ident{axioms}} \derives G$ por el puente entre los dos cálculos
  (capítulo~\ref{cap:kernel}).
\item El punto fijo, leído de izquierda a derecha, dice
  $G \Rightarrow \lnot\Prov(\codigo{G})$. Con (2), la teoría demuestra
  $\lnot\Prov(\codigo{G})$.
\item Junto a (1): la teoría demuestra una fórmula y su negación. Luego demuestra $\bot$.
\end{enumerate}
Si la teoría es consistente, (4) es imposible. Luego $G$ no es demostrable. $\square$
\end{matematica}

Obsérvese que no se ha hablado ni una vez de si $G$ es \emph{verdadera}. El argumento entero es
sintáctico: entra una demostración hipotética y sale una contradicción.

\section{El mismo argumento, en Lean}

Aquí está, entero, sin recortes. Es uno de los pocos sitios de este libro donde merece la pena
imprimir la prueba completa y no sólo el enunciado, porque las cinco líneas del cuerpo son las
cuatro del argumento de arriba más la contradicción:

\leanfrag{goedel_first_unprovable}

\selee{«si la teoría es consistente y $G$ es un punto fijo de $\lnot\Prov$, entonces $G$ no tiene
demostración finitaria»}

Léase la correspondencia línea a línea: \ident{repr_pos'} es D1; \ident{prf_to_derives} es el
puente entre los cálculos; \ident{and_elim_left} extrae la dirección del bicondicional que hace
falta; y la hipótesis \texttt{hcon} —la consistencia— es lo único que se pone de fuera.

\begin{leccion}[Por qué el teorema está partido en dos declaraciones]
Nótese que este enunciado no habla de \emph{la} sentencia de Gödel: habla de \textbf{cualquier}
punto fijo $G$. Es el argumento, aislado de la construcción que lo alimenta.

Eso no es elegancia gratuita. Cuando la construcción del punto fijo resultó estar apoyada en un
axioma inconsistente y hubo que rehacerla entera (capítulo~\ref{cap:numerales}), \textbf{este
teorema no se tocó}: bastó con enchufarle el punto fijo nuevo. Separar el argumento de su insumo
es lo que hizo que una reparación de los cimientos no obligara a rehacer el tejado.
\end{leccion}

Y la descarga, que es una línea:

\leanfrag{goedel_first_numeral}

\begin{observacion}
\emph{Cero postulados gödelianos.} No hay ninguna hipótesis del tipo «supongamos que $\Prov$
representa fielmente la demostrabilidad», que es donde las formalizaciones ingenuas esconden el
trabajo. Lo único que entra es la consistencia. Todo lo demás son teoremas de las partes
anteriores.
\end{observacion}

\section{Qué consistencia, exactamente}

La hipótesis se llama \ident{ConsistentOmega}, y conviene mirarla porque su nombre engaña:

\leanfrag{ConsistentOmega}

\selee{«la teoría no demuestra $\bot$ en el cálculo con la $\omega$-regla»}

Es \textbf{consistencia simple} —la más débil que se puede pedir—, enunciada sobre el cálculo
fuerte $\derives$. Y de ahí se hereda para el cálculo finitario:

\leanfrag{consistentH_of_omega}

\begin{muro}[Dos nombres a una palabra de distancia, y no son lo mismo]
En este proyecto conviven \ident{ConsistentOmega} y \ident{OmegaConsistent}. Se leen casi igual y
son propiedades distintas — y la que \emph{suena} más fuerte es la más débil:

\leanfrag{OmegaConsistent}

\ident{ConsistentOmega} es «el cálculo $\omega$ es consistente»: no demuestra $\bot$.
\ident{OmegaConsistent} es la \textbf{$\omega$-consistencia} de verdad: que la teoría no demuestre
$\exists x.\,A(x)$ mientras refuta $A(\bar n)$ para todos los testigos estándar.

La primera es la hipótesis de este capítulo. La segunda hace falta para la otra mitad. Confundirlas
en una lectura rápida convierte un teorema honesto en una exageración, y el orden de las palabras
es lo único que las separa.
\end{muro}

\section{Sobre qué teoría se ha demostrado esto}

\avisoreparacion

Conviene ser preciso sobre qué significa eso aquí. \ident{goedel_first_numeral} es un teorema
\textbf{real} sobre la teoría \textbf{reparada}: se retiró el axioma que hacía derivable $\bot$, y
lo que queda no tiene ninguna inconsistencia conocida. Pero «ninguna conocida» no es «ninguna», y
este libro no va a decir lo segundo.

\begin{leccion}[Y la hipótesis se lleva bien con esa cautela]
Hay una elegancia real en que el teorema tome la consistencia como \emph{hipótesis} en vez de
darla por buena. El enunciado no afirma que la teoría sea consistente: afirma que
\textbf{si} lo es, $G$ no es demostrable. Esa forma condicional es la que sobrevive a un
descubrimiento como el de julio de 2026 — y, de hecho, es la única forma en la que el segundo
teorema de Gödel permite enunciar nada sobre la consistencia de un sistema desde dentro.
\end{leccion}

\section{La mitad que falta}

\avisomediagodel

El teorema clásico dice que $G$ es \textbf{indecidible}: ni demostrable ni refutable. Lo que este
proyecto tiene demostrado es la primera mitad. La segunda existe como enunciado y arrastra una
hipótesis que hoy no está descargada — es el asunto del capítulo siguiente.

La razón de fondo no es un descuido, y ya estaba en el teorema original.

\begin{leccion}[Las dos mitades nunca costaron lo mismo]
Gödel demostró $\noderives G$ con la consistencia simple, y necesitó la
\textbf{$\omega$-consistencia} para $\noderives \lnot G$. La asimetría es del argumento, no de la
formalización: la primera mitad sólo tiene que impedir que la teoría demuestre algo y su negación;
la segunda tiene que impedir que la teoría se equivoque \emph{sobre los testigos}, y eso es una
exigencia más fuerte.

Rosser la eliminó en 1936 —cambiando la sentencia, no el argumento: su $R$ dice «si hay una
demostración mía, hay antes una demostración de mi negación»— y consiguió las dos mitades con
consistencia simple.\footnote{J.~B. Rosser, «Extensions of some theorems of Gödel and Church»,
\emph{The Journal of Symbolic Logic} 1(3), 1936, pp.~87--91.} Este proyecto \textbf{no} tomó esa
ruta: construyó la sentencia de Gödel y dejó la segunda mitad como una hipótesis explícita. Es una
decisión, y decirla es parte de decir qué se ha demostrado.
\end{leccion}

\section{Lo que queda dicho}

\begin{center}\small
\begin{tabularx}{\linewidth}{@{}l >{\raggedright\arraybackslash}X l@{}}
\toprule
mitad & qué hace falta & estado \\
\midrule
$\noderives G$ & consistencia simple, y nada más & \textbf{demostrado} \\
$\noderives \lnot G$ & reflexión —o $\omega$-consistencia más un verificador negativo— & hipótesis sin descargar \\
\bottomrule
\end{tabularx}
\end{center}

Media indecidibilidad no es un resultado a medias: $\noderives G$ es, por sí solo, el enunciado que
hunde el programa de Hilbert. Hay una sentencia aritmética que la teoría no puede demostrar, y su
número de Gödel se puede calcular. Lo que falta es saber si la teoría tampoco puede refutarla — y
en eso, como se verá, el proyecto tiene el enunciado escrito y ni una línea de la construcción.
